\section{AMC 8 1993}
        \begin{enumerate}
        	\item (\textit{AMC 8 1993}) Which pair of numbers does NOT have a product equal to $36$?


 $\text{(A)}\ \{ -4,-9\} \qquad \text{(B)}\ \{ -3,-12\} \qquad \text{(C)}\ \left\{ \dfrac{1}{2},-72\right\} \qquad \text{(D)}\ \{ 1,36\} \qquad \text{(E)}\ \left\{\dfrac{3}{2},24\right\}$


 

	\item (\textit{AMC 8 1993}) When the fraction $\dfrac{49}{84}$ is expressed in simplest form, then the sum of the numerator and the denominator will be 


 $\text{(A)}\ 11 \qquad \text{(B)}\ 17 \qquad \text{(C)}\ 19 \qquad \text{(D)}\ 33 \qquad \text{(E)}\ 133$


 

	\item (\textit{AMC 8 1993}) Which of the following numbers has the largest prime factor?


 $\text{(A)}\ 39 \qquad \text{(B)}\ 51 \qquad \text{(C)}\ 77 \qquad \text{(D)}\ 91 \qquad \text{(E)}\ 121$


 

	\item (\textit{AMC 8 1993}) $1000\times 1993 \times 0.1993 \times 10 =$


 $\text{(A)}\ 1.993\times 10^3 \qquad \text{(B)}\ 1993.1993 \qquad \text{(C)}\ (199.3)^2 \qquad \text{(D)}\ 1,993,001.993 \qquad \text{(E)}\ (1993)^2$


 

	\item (\textit{AMC 8 1993}) Which one of the following bar graphs could represent the data from the circle graph?


 
                    \begin{figure}[H]
                        \centering
                        \includegraphics[width=0.5\textwidth]{Problem/AMC_8_AJHSME_1993_5_Problem_Figure_1.png}
                    \end{figure}
                    


 
                    \begin{figure}[H]
                        \centering
                        \includegraphics[width=0.5\textwidth]{Problem/AMC_8_AJHSME_1993_5_Problem_Figure_2.png}
                    \end{figure}
                    


 

	\item (\textit{AMC 8 1993}) A can of soup can feed $3$ adults or $5$ children.  If there are $5$ cans of soup and $15$ children are fed, then how many adults would the remaining soup feed?


 $\text{(A)}\ 5 \qquad \text{(B)}\ 6 \qquad \text{(C)}\ 7 \qquad \text{(D)}\ 8 \qquad \text{(E)}\ 25$


 

	\item (\textit{AMC 8 1993}) $3^3+3^3+3^3 =$


 $\text{(A)}\ 3^4 \qquad \text{(B)}\ 9^3 \qquad \text{(C)}\ 3^9 \qquad \text{(D)}\ 27^3 \qquad \text{(E)}\ 3^{27}$


 

	\item (\textit{AMC 8 1993}) To control her blood pressure, Jill's grandmother takes one half of a pill every other day.  If one supply of medicine contains $60$ pills, then the supply of medicine would last approximately


 $\text{(A)}\ 1\text{ month} \qquad \text{(B)}\ 4\text{ months} \qquad \text{(C)}\ 6\text{ months} \qquad \text{(D)}\ 8\text{ months} \qquad \text{(E)}\ 1\text{ year}$


 

	\item (\textit{AMC 8 1993}) Consider the operation $*$ defined by the following table:


 
 \[\begin{tabular}{c|cccc} * & 1 & 2 & 3 & 4 \\ \hline 1 & 1 & 2 & 3 & 4 \\ 2 & 2 & 4 & 1 & 3 \\ 3 & 3 & 1 & 4 & 2 \\ 4 & 4 & 3 & 2 & 1 \end{tabular}\]


 For example, $3*2=1$.  Then $(2*4)*(1*3)=$


 $\text{(A)}\ 1 \qquad \text{(B)}\ 2 \qquad \text{(C)}\ 3 \qquad \text{(D)}\ 4 \qquad \text{(E)}\ 5$


 

	\item (\textit{AMC 8 1993}) This line graph represents the price of a trading card during the first $6$ months of $1993$.


 
                    \begin{figure}[H]
                        \centering
                        \includegraphics[width=0.5\textwidth]{Problem/AMC_8_AJHSME_1993_10_Problem_Figure_1.png}
                    \end{figure}
                    


 The greatest monthly drop in price occurred during


 $\text{(A)}\ \text{January} \qquad \text{(B)}\ \text{March} \qquad \text{(C)}\ \text{April} \qquad \text{(D)}\ \text{May} \qquad \text{(E)}\ \text{June}$


 

	\item (\textit{AMC 8 1993}) Consider this histogram of the scores for $81$ students taking a test:


 
                    \begin{figure}[H]
                        \centering
                        \includegraphics[width=0.5\textwidth]{Problem/AMC_8_AJHSME_1993_11_Problem_Figure_1.png}
                    \end{figure}
                    


 The median is in the interval labeled 


 $\text{(A)}\ 60 \qquad \text{(B)}\ 65 \qquad \text{(C)}\ 70 \qquad \text{(D)}\ 75 \qquad \text{(E)}\ 80$


 

	\item (\textit{AMC 8 1993}) If each of the three operation signs, $+$, $-$, $\times$, is used exactly ONCE in one of the blanks in the expression


 
 \[5\hspace{1 mm}\underline{\hspace{4 mm}}\hspace{1 mm}4\hspace{1 mm}\underline{\hspace{4 mm}}\hspace{1 mm}6\hspace{1 mm}\underline{\hspace{4 mm}}\hspace{1 mm}3\]


 then the value of the result could equal


 $\text{(A)}\ 9 \qquad \text{(B)}\ 10 \qquad \text{(C)}\ 15 \qquad \text{(D)}\ 16 \qquad \text{(E)}\ 19$


 

	\item (\textit{AMC 8 1993}) The word "\textbf{HELP}" in block letters is painted in black with strokes $1$ unit wide on a $5$ by $15$ rectangular white sign with dimensions as shown.  The area of the white portion of the sign, in square units, is


 
                    \begin{figure}[H]
                        \centering
                        
\includegraphics[width=0.5\textwidth]{Problem/AMC_8_AJHSME_1993_13_Problem_Figure_1.png}
                    \end{figure}
                    


 $\text{(A)}\ 30 \qquad \text{(B)}\ 32 \qquad \text{(C)}\ 34 \qquad \text{(D)}\ 36 \qquad \text{(E)}\ 38$


 

	\item (\textit{AMC 8 1993}) The nine squares in the table shown are to be filled so that every row and every column contains each of the numbers $1,2,3$.  Then $A+B=$


 
 \[\begin{tabular}{|c|c|c|} \hline 1 & & \\ \hline  & 2 & A \\ \hline  & & B \\ \hline \end{tabular}\]


 $\text{(A)}\ 2 \qquad \text{(B)}\ 3 \qquad \text{(C)}\ 4 \qquad \text{(D)}\ 5 \qquad \text{(E)}\ 6$


 

	\item (\textit{AMC 8 1993}) The arithmetic mean (average) of four numbers is $85$.  If the largest of these numbers is $97$, then the mean of the remaining three numbers is 


 $\text{(A)}\ 81.0 \qquad \text{(B)}\ 82.7 \qquad \text{(C)}\ 83.0 \qquad \text{(D)}\ 84.0 \qquad \text{(E)}\ 84.3$


 

	\item (\textit{AMC 8 1993}) $\dfrac{1}{1+\dfrac{1}{2+\dfrac{1}{3}}} =$


 $\text{(A)}\ \dfrac{1}{6} \qquad \text{(B)}\ \dfrac{3}{10} \qquad \text{(C)}\ \dfrac{7}{10} \qquad \text{(D)}\ \dfrac{5}{6} \qquad \text{(E)}\ \dfrac{10}{3}$


 

	\item (\textit{AMC 8 1993}) Square corners, $5$ units on a side, are removed from a $20$ unit by $30$ unit rectangular sheet of cardboard.  The sides are then folded to form an open box.  The surface area, in square units, of the interior of the box is


 
                    \begin{figure}[H]
                        \centering
                        \includegraphics[width=0.5\textwidth]{Problem/AMC_8_AJHSME_1993_17_Problem_Figure_1.png}
                    \end{figure}
                    


 $\text{(A)}\ 300 \qquad \text{(B)}\ 500 \qquad \text{(C)}\ 550 \qquad \text{(D)}\ 600 \qquad \text{(E)}\ 1000$


 

	\item (\textit{AMC 8 1993}) The rectangle shown has length $AC=32$, width $AE=20$, and $B$ and $F$ are midpoints of $\overline{AC}$ and $\overline{AE}$, respectively.  The area of quadrilateral $ABDF$ is


 
                    \begin{figure}[H]
                        \centering
                        \includegraphics[width=0.5\textwidth]{Problem/AMC_8_AJHSME_1993_18_Problem_Figure_1.png}
                    \end{figure}
                    


 $\text{(A)}\ 320 \qquad \text{(B)}\ 325 \qquad \text{(C)}\ 330 \qquad \text{(D)}\ 335 \qquad \text{(E)}\ 340$


 

	\item (\textit{AMC 8 1993}) $(1901+1902+1903+\cdots + 1993) - (101+102+103+\cdots + 193) =$


 $\text{(A)}\ 167,400 \qquad \text{(B)}\ 172,050 \qquad \text{(C)}\ 181,071 \qquad \text{(D)}\ 199,300 \qquad \text{(E)}\ 362,142$


 

	\item (\textit{AMC 8 1993}) When $10^{93}-93$ is expressed as a single whole number, the sum of the digits is


 $\text{(A)}\ 10 \qquad \text{(B)}\ 93 \qquad \text{(C)}\ 819 \qquad \text{(D)}\ 826 \qquad \text{(E)}\ 833$


 

	\item (\textit{AMC 8 1993}) If the length of a rectangle is increased by $20\%$ and its width is increased by $50\%$, then the area is increased by 


 $\text{(A)}\ 10\% \qquad \text{(B)}\ 30\% \qquad \text{(C)}\ 70\% \qquad \text{(D)}\ 80\% \qquad \text{(E)}\ 100\%$


 

	\item (\textit{AMC 8 1993}) Pat Peano has plenty of 0's, 1's, 3's, 4's, 5's, 6's, 7's, 8's and 9's, but he has only twenty-two 2's.  How far can he number the pages of his scrapbook with these digits?


 $\text{(A)}\ 22 \qquad \text{(B)}\ 99 \qquad \text{(C)}\ 112 \qquad \text{(D)}\ 119 \qquad \text{(E)}\ 199$


 

	\item (\textit{AMC 8 1993}) Five runners, $P$, $Q$, $R$, $S$, $T$, have a race, and $P$ beats $Q$, $P$ beats $R$, $Q$ beats $S$, and $T$ finishes after $P$ and before $Q$.  Who could NOT have finished third in the race?


 $\text{(A)}\ P\text{ and }Q \qquad \text{(B)}\ P\text{ and }R \qquad \text{(C)}\ P\text{ and }S \qquad \text{(D)}\ P\text{ and }T \qquad \text{(E)}\ P,S\text{ and }T$


 

	\item (\textit{AMC 8 1993}) What number is directly above $142$ in this array of numbers?


 
 \[\begin{array}{cccccc}  & & & 1 & & \\ & & 2 & 3 & 4 & \\ & 5 & 6 & 7 & 8 & 9 \\ 10 & 11 & 12 & \cdots & & \\ \end{array}\]


 $\text{(A)}\ 99 \qquad \text{(B)}\ 119 \qquad \text{(C)}\ 120 \qquad \text{(D)}\ 121 \qquad \text{(E)}\ 122$


 

	\item (\textit{AMC 8 1993}) A checkerboard consists of one-inch squares.  A square card, $1.5$ inches on a side, is placed on the board so that it covers part or all of the area of each of $n$ squares.  The maximum possible value of $n$ is


 $\text{(A)}\ 4\text{ or }5 \qquad \text{(B)}\ 6\text{ or }7\qquad \text{(C)}\ 8\text{ or }9 \qquad \text{(D)}\ 10\text{ or }11 \qquad \text{(E)}\ 12\text{ or more}$


 

\end{enumerate}